% Part: many-valued-logic
% Chapter: sequent-calculus
% Section: propositional-rules

\documentclass[../../../include/open-logic-section]{subfiles}

\begin{document}

\olfileid{mvl}{seq}{prl}

\olsection{Aturan Proposisional kanggo Logika Tartamtu}

Aturan inferensi kanggo sawijining konektif ing kalkulus
sekuen $n$ sisi mung gumantung marang fungsi kabeneran khas
konektif kasebut. Mula, yen sawijining konektif ditetepake
nganggo fungsi kabeneran kang padha ing logika-logika kang beda,
aturan sekuen $n$ sisi kanggo konektif iku uga padha ing
logika-logika kasebut.

\subsection{Aturan kanggo $\lnot$}

Aturan kanggo $\lnot$ ing ngisor iki lumaku ing logika
\L ukasiewicz lan Kleene, uga variasine.

\begin{defish}
  \begin{center}
\AxiomC{$ \Gamma \nSequent \Pi \nSequent \Delta, !A $}
\RightLabel{\iR{\lnot}{\False}}
\UnaryInfC{$ \lnot !A, \Gamma \nSequent \Pi \nSequent \Delta$}
\DisplayProof
\\[2ex]
\AxiomC{$ \Gamma \nSequent !A, \Pi \nSequent \Delta $}
\RightLabel{\iR{\lnot}{\Undef}}
\UnaryInfC{$\Gamma \nSequent \lnot !A, \Pi \nSequent \Delta$}
\DisplayProof
\\[2ex]
\AxiomC{$!A, \Gamma \nSequent \Pi \nSequent \Delta$}
\RightLabel{\iR{\lnot}{\True}}
\UnaryInfC{$\Gamma \nSequent \Pi \nSequent \Delta,  \lnot !A$}
\DisplayProof
  \end{center}
\end{defish}

Aturan kanggo $\lnot$ ing ngisor iki lumaku ing
logika G\"odel.

\begin{defish}
\AxiomC{$ \Gamma \nSequent !A, \Pi \nSequent \Delta, !A $}
\RightLabel{\iR{\lnot}{\False}[\LogGod]}
\UnaryInfC{$ \lnot !A, \Gamma \nSequent \Pi \nSequent \Delta$}
\DisplayProof
\hfill
\AxiomC{$!A, \Gamma \nSequent \Pi \nSequent \Delta$}
\RightLabel{\iR{\lnot}{\True}[\LogGod]}
\UnaryInfC{$\Gamma \nSequent \Pi \nSequent \Delta,  \lnot !A$}
\DisplayProof
\hfill
\end{defish}

(Ing logika G\"odel, $\lnot !A$ ora bisa nduweni
nilai~$\Undef$, mula ora ana aturan kanggo panggonan tengah.)

\subsection{Aturan kanggo $\land$}

Iki aturan kanggo $\land$ ing logika \L ukasiewicz,
Kleene kang kuwat, lan G\"odel.

\begin{defish}
\begin{center}
  \AxiomC{$!A, !B, \Gamma \nSequent \Pi \nSequent \Delta$}
  \RightLabel{$\iR{\land}{\False}$}
  \UnaryInfC{$!A \land !B, \Gamma \nSequent \Pi \nSequent \Delta$}
  \DisplayProof
\\[2ex]
  \AxiomC{$\Gamma \nSequent !A, \Pi \nSequent !A, \Delta$}
  \AxiomC{$\Gamma \nSequent !B, \Pi \nSequent !B, \Delta$}
  \AxiomC{$\Gamma \nSequent !A, !B, \Pi \nSequent \Delta$}
  \RightLabel{$\iR\land\Undef$}
  \TrinaryInfC{$\Gamma \nSequent !A \land !B, \Pi \nSequent \Delta$}
  \DisplayProof
  \\[2ex]  
  \AxiomC{$\Gamma \nSequent \Pi \nSequent \Delta, !A$}
  \AxiomC{$\Gamma \nSequent \Pi \nSequent \Delta, !B$}
  \RightLabel{$\iR\land\True$}
  \BinaryInfC{$\Gamma \nSequent \Pi \nSequent \Delta, !A \land !B$}
  \DisplayProof
\end{center}
\end{defish}

\subsection{Aturan kanggo $\lor$}

Iki aturan kanggo $\lor$ ing logika \L ukasiewicz,
Kleene kang kuwat, lan G\"odel.

\begin{defish}
\begin{center}
 \AxiomC{$!A, \Gamma \nSequent \Pi \nSequent \Delta$}
  \AxiomC{$!B, \Gamma \nSequent \Pi \nSequent \Delta$}
  \RightLabel{$\iR\lor\False$}
  \BinaryInfC{$!A \lor !B, \Gamma \nSequent \Pi \nSequent \Delta$}
  \DisplayProof
\\[2ex]
  \AxiomC{$!A, \Gamma \nSequent !A, \Pi \nSequent \Delta$}
  \AxiomC{$!B, \Gamma \nSequent !B, \Pi \nSequent \Delta$}
  \AxiomC{$\Gamma \nSequent !A, !B, \Pi \nSequent \Delta$}
  \RightLabel{$\iR\lor\Undef$}
  \TrinaryInfC{$\Gamma \nSequent !A \lor !B, \Pi \nSequent \Delta$}
  \DisplayProof
  \\[2ex]  
  \AxiomC{$\Gamma \nSequent \Pi \nSequent \Delta, !A, !B$}
  \RightLabel{$\iR\lor\True$}
  \UnaryInfC{$\Gamma \nSequent \Pi \nSequent \Delta, !A \lor !B$}
  \DisplayProof
\end{center}
\end{defish}

\subsection{Aturan kanggo $\lif$}

Iki aturan kanggo $\lif$ ing logika \L ukasiewicz.

\begin{defish}
\begin{center}
  \AxiomC{$\Gamma \nSequent \Pi \nSequent \Delta, !A$}
  \AxiomC{$!B, \Gamma \nSequent \Pi \nSequent \Delta$}
  \RightLabel{$\iR\lif\False[\LogLuk[3]]$}
  \BinaryInfC{$!A \lif !B, \Gamma \nSequent \Pi \nSequent \Delta$}
  \DisplayProof
  \\[2ex]
  \AxiomC{$\Gamma \nSequent !A, !B, \Pi \nSequent \Delta$}
  \AxiomC{$!B, \Gamma \nSequent \Pi \nSequent \Delta, !A$}
  \RightLabel{$\iR\lif\Undef[\LogLuk[3]]$}
  \BinaryInfC{$\Gamma \nSequent !A \lif !B, \Pi \nSequent \Delta$}
  \DisplayProof
\\[2ex]
  \AxiomC{$!A, \Gamma \nSequent !B, \Pi \nSequent \Delta, !B$}
  \AxiomC{$!A, \Gamma \nSequent !A, \Pi \nSequent \Delta, !B$}
  \RightLabel{$\iR{\lif}{\True}[\LogLuk[3]]$}
  \BinaryInfC{$\Gamma \nSequent \Pi \nSequent \Delta, !A \lif !B$}
  \DisplayProof
\end{center}
\end{defish}

Iki aturan kanggo $\lif$ ing logika Kleene kang kuwat.

\begin{defish}
\begin{center}
  \AxiomC{$\Gamma \nSequent \Pi \nSequent \Delta, !A$}
  \AxiomC{$!B, \Gamma \nSequent \Pi \nSequent \Delta$}
  \RightLabel{$\iR\lif\False[\LogKs]$}
  \BinaryInfC{$!A \lif !B, \Gamma \nSequent \Pi \nSequent \Delta$}
  \DisplayProof
  \\[2ex]
  \AxiomC{$!B, \Gamma \nSequent !B, \Pi \nSequent \Delta$}
  \AxiomC{$\Gamma \nSequent !A, !B, \Pi \nSequent \Delta$}
  \AxiomC{$\Gamma \nSequent !A, \Pi \nSequent \Delta, !A$}
  \RightLabel{$\iR\lif\Undef[\LogKs]$}
  \TrinaryInfC{$\Gamma \nSequent !A \lif !B, \Pi \nSequent \Delta$}
  \DisplayProof
\\[2ex]
  \AxiomC{$!A, \Gamma \nSequent \Pi \nSequent \Delta, !B$}
  \RightLabel{$\iR{\lif}{\True}[\LogKs]$}
  \UnaryInfC{$\Gamma \nSequent \Pi \nSequent \Delta, !A \lif !B$}
  \DisplayProof
\end{center}
\end{defish}

Iki aturan kanggo $\lif$ ing logika G\"odel.

\begin{defish}
\begin{center}
  \AxiomC{$\Gamma \nSequent !A, \Pi \nSequent \Delta, !A$}
  \AxiomC{$!B, \Gamma \nSequent \Pi \nSequent \Delta$}
  \RightLabel{$\iR\lif\False[\LogGod[3]]$}
  \BinaryInfC{$!A \lif !B, \Gamma \nSequent \Pi \nSequent \Delta$}
  \DisplayProof
  \\[2ex]
  \AxiomC{$\Gamma \nSequent !B, \Pi \nSequent \Delta$}
  \AxiomC{$\Gamma \nSequent \Pi \nSequent \Delta, !A$}
  \RightLabel{$\iR\lif\Undef[\LogGod[3]]$}
  \BinaryInfC{$\Gamma \nSequent !A \lif !B, \Pi \nSequent \Delta$}
  \DisplayProof
\\[2ex]
  \AxiomC{$!A, \Gamma \nSequent !B, \Pi \nSequent \Delta, !B$}
  \AxiomC{$!A, \Gamma \nSequent !A, \Pi \nSequent \Delta, !B$}
  \RightLabel{$\iR{\lif}{\True}[\LogGod[3]]$}
  \BinaryInfC{$\Gamma \nSequent \Pi \nSequent \Delta, !A \lif !B$}
  \DisplayProof
\end{center}
\end{defish}

\begin{sidewaysfigure}
  \begin{center}  
  \AxiomC{$A \nSequent A \nSequent A$}
  \RightLabel{\iR \Weakening \True}
  \UnaryInfC{$A \nSequent A \nSequent B, A$}
  \RightLabel{\iR \Weakening \Undef}
  \UnaryInfC{$A \nSequent B, A \nSequent B, A$}
  \RightLabel{\iR \Weakening \Undef}
  \UnaryInfC{$A \nSequent A, B, A \nSequent B, A$}
  \AxiomC{$A \nSequent A \nSequent A$}
  \RightLabel{\iR \Weakening \True}
  \UnaryInfC{$A \nSequent A \nSequent A, A$}
  \RightLabel{\iR \Weakening \True}
  \UnaryInfC{$A \nSequent A \nSequent B, A, A$}
  \RightLabel{\iR \Weakening \False}
  \UnaryInfC{$B, A \nSequent A \nSequent B, A, A$}
  \RightLabel{$\iR\lif\Undef$}
  \BinaryInfC{$A \nSequent A \lif B, A \nSequent B, A$}
  \AxiomC{$B \nSequent B \nSequent B$}
  \RightLabel{\iR \Weakening \Undef}
  \UnaryInfC{$B \nSequent A, B \nSequent B$}
  \RightLabel{\iR \Exchange \Undef}
  \UnaryInfC{$B \nSequent B, A \nSequent B$}
  \RightLabel{\iR \Weakening \Undef}
  \UnaryInfC{$B \nSequent A, B, A \nSequent B$}
  \RightLabel{\iR \Weakening \False}
  \UnaryInfC{$A, B \nSequent A, B, A \nSequent B$}
  \RightLabel{\iR \Exchange \False}
  \UnaryInfC{$B, A \nSequent A, B, A \nSequent B$}
  \AxiomC{$A \nSequent A \nSequent A$}
  \RightLabel{\iR \Weakening \True}
  \UnaryInfC{$A \nSequent A \nSequent B, A$}
  \RightLabel{\iR \Weakening \False}
  \UnaryInfC{$B, A \nSequent A \nSequent B, A$}
  \RightLabel{\iR \Weakening \False}
  \UnaryInfC{$B, B, A \nSequent A \nSequent B, A$}
  \RightLabel{$\iR\lif\Undef$}
  \BinaryInfC{$B, A \nSequent A \lif B, A \nSequent B$}
  \RightLabel{$\iR\lif\False$}
  \BinaryInfC{$A \lif B, A \nSequent A \lif B, A \nSequent B$}
  \DisplayProof
  \end{center}
  \caption{Tuladha !!{derivation} ing~$\LogLuk[3]$}
\end{sidewaysfigure}

\end{document}
